\documentclass{article}

\usepackage{mathtools}
\usepackage{amssymb}

\DeclareMathOperator{\sinc}{sinc}

\title{Breviar teoretic -- Matematici Speciale}
\author{Andrei Florea}
\date{19 martie 2026}

\begin{document}
\maketitle

\section{Analiză Fourier}
\subsection{Serii Fourier}

Seria Fourier asociată funcției $f$ (funcție periodică cu perioada $T=2\pi$):
\begin{displaymath}
    \mathrm{S}_f(t)=\frac{a_0}{2}+\sum_{n\geq1}^\infty (a_n\cos nt+b_n\sin nt)
\end{displaymath}

unde $a_n$ și $b_n$ sunt coeficienții Fourier asociați funcției $f$:
\begin{gather*}
    a_n=\frac{1}{\pi}\int_{-\pi}^\pi{f(t)\cos nt\,\mathrm{d}t}\quad(n\geq0) \\
    b_n=\frac{1}{\pi}\int_{-\pi}^\pi{f(t)\sin nt\,\mathrm{d}t}\quad(n\geq1)
\end{gather*}
% Orice funcție periodică se poate transforma într-o funcție periodică cu perioada $2\pi$:
% \begin{displaymath}
%     g(t)=f(t\cdot\frac{T}{2\pi})
% \end{displaymath}
Observație pentru funcțiile a căror formă nu este cunoscută pe $(-\pi,\pi)$:
\begin{displaymath}
    \int_a^{a+T}{f(t)\,\mathrm{d}t}=\int_0^T{f(t)\,\mathrm{d}t},\,\forall a\in\mathbb{R}
\end{displaymath}
Formula lui Parseval:
\begin{displaymath}
    \frac{{|a_0|}^2}{2}+\sum_{n\geq1}^\infty ({|a_n|}^2+{|b_n|}^2) = \frac{1}{\pi}\int_{-\pi}^{\pi}{|f(t)|^2\,\mathrm{d}t}
\end{displaymath}

\subsubsection{Informații utile}

\begin{enumerate}
\item
\begin{enumerate}
\item \(\displaystyle{f(-x)=-f(x)\Longrightarrow\int_{-a}^a{f(x)\,\mathrm{d}x}=0}\)
\item \(\displaystyle{f(-x)=f(x)\Longrightarrow\int_{-a}^a{f(x)\,\mathrm{d}x}=2\int_0^a{f(x)\,\mathrm{d}x}}\)
\end{enumerate}

\item
\begin{enumerate}
\item Produsul a două funcții de aceeași paritate este o funcție pară.
\item Produsul a două funcții de parități diferite este o funcție impară.
\end{enumerate}
\item
\begin{enumerate}
\item Prin compunerea unei funcții pare cu o funcție de orice paritate rezultă o funcție pară.
\item Prin compunerea a două funcții impare rezultă o funcție impară.
\end{enumerate}

\item
\begin{enumerate}
\item \(\displaystyle{\int{\cos nt\,\mathrm{d}t}=\frac{\sin nt}{n}+C}\)
\item \(\displaystyle{\int{\sin nt\,\mathrm{d}t}=-\frac{\cos nt}{n}+C}\)
\end{enumerate}

\item
\begin{enumerate}
\item \(\sin n\pi=0\)
\item \(\cos n\pi=(-1)^n\)
\end{enumerate}

\item Funcția sinus atenuat
\begin{displaymath}
    \sinc\colon\mathbb{R}\to\mathbb{R},\quad
    \sinc x=
    \begin{cases}
        \dfrac{\sin x}{x}, & x \ne 0 \\
        1, & x=0
    \end{cases}
\end{displaymath}
\end{enumerate}

\subsection{Transformarea Fourier}

O funcție $f\colon\mathbb{R}\to\mathbb{C}$ este absolut integrabilă ($f \in L^1$) dacă
\[\int_{-\infty}^{\infty}{|f(t)|\,\mathrm{d}t}<\infty\]
Transformata Fourier a unui semnal $f$ este $F\colon\mathbb{R}\to\mathbb{C}$:
\[F(\omega)=\int_{-\infty}^{\infty}{f(t)e^{-i\omega t}\,\mathrm{d}t}\]
Notăm \(f(t)\overset{\mathcal{F}}{\longrightarrow}F(\omega)\). \\
Transformata Fourier inversă a funcției $F$ este $f\colon\mathbb{R}\to\mathbb{C}$:
\[
    \frac{1}{2\pi}\int_{-\omega}^{\omega}{F(\omega)e^{i\omega t}\,\mathrm{d}t}=
    \begin{cases}
        f(t), & \text{f continuă în t} \\
        \dfrac{f(t-0)+f(t+0)}{2}, & \text{în rest}
    \end{cases}
\]
Condiții:
\begin{enumerate}
\item $f\in L^1,F\in L^1$, $f$ continuă; sau
\item $f\in L^1$ și $f$ îndeplinește condițiile lui Dirichlet
\end{enumerate}

\subsection{Transformarea Fourier Discretă}

\end{document}
